% =============================================================================
\section{Simulations et scénarios réalistes}
% =============================================================================

\begin{frame}{Simulation : charges historiques et ressources supposées}
  \small
  \begin{columns}[T,onlytextwidth]
    \begin{column}{0.32\textwidth}
      \begin{block}{Données}
        \centering
        {\Large\bfseries 14\,434}\\
        cas exploitables\\[5pt]
        73 exclusions\\[5pt]
        2019--2021 : estimation\\
        2022 : exécution cachée
      \end{block}
    \end{column}
    \begin{column}{0.32\textwidth}
      \begin{block}{Deux charges}
        \centering
        {\Large\bfseries 58} cas / 5 jours\\[5pt]
        {\Large\bfseries 279} cas / 22 jours\\[5pt]
        Cohorte test 2022\\
        Patients hors entraînement
      \end{block}
    \end{column}
    \begin{column}{0.32\textwidth}
      \begin{block}{Scénarios}
        \centering
        2 salles, 08h--17h\\
        15 min de remise en état\\[5pt]
        42 ou 10 lits \textbf{synthétiques}\\
        Salle 1 fermée de 10h à 12h
      \end{block}
    \end{column}
  \end{columns}
  \vspace{0.23cm}
  \begin{exampleblock}{Comparaison appariée}
    \small Mêmes arrivées, mêmes durées réellement observées, mêmes
    ressources. Trois informations : médianes, prédictions ML, oracle.
  \end{exampleblock}
  \source{Fenêtres du 3 au 7 et du 3 au 31 janvier 2022 ; charges de travail, pas preuve de passage à l'échelle des solveurs.}
\end{frame}

\begin{frame}{Intégration SMA : planifier, simuler, corriger}
  \small
  \centering
  \vspace{0.25cm}
  \begin{tikzpicture}[
      card/.style={draw=McmTerracotta,rounded corners=2pt,fill=white,
                   minimum height=2.05cm,text width=3.10cm,align=center,
                   inner sep=5pt,font=\small},
      flow/.style={-{Latex[length=2.5mm]},thick,draw=McmTerracotta}]
    \node[card] (input) at (0,0) {\textbf{Entrées disponibles}\\[3pt]
      Interventions à programmer\\Durées prévues : salle et séjour restant};
    \node[card] (solver) at (4.15,0) {\textbf{Planificateur}\\[3pt]
      Glouton ou métaheuristique\\Planning proposé puis validé};
    \node[card] (mesa) at (8.30,0) {\textbf{Exécution Mesa}\\[3pt]
      Agents-épisodes\\Salles à la minute, lits sur plusieurs jours};
    \draw[flow] (input) -- (solver);
    \draw[flow] (solver) -- (mesa);
    \draw[flow] (mesa.south) -- ++(0,-0.80) -| (solver.south);
    \node[font=\scriptsize,text=McmMuted,fill=McmWarm,inner sep=2pt]
      at (6.22,-1.53) {Observations et replanification des cas en attente};
  \end{tikzpicture}
  \vspace{0.23cm}

  \begin{exampleblock}{Le rôle du SMA ici}
    \small Le solveur utilise des durées estimées. Mesa révèle les résultats
    réels au coordinateur.
  \end{exampleblock}
  \source{Implémenté : hospital\_sim (Mesa 3.5.1). Les lits sont des ressources ; l'ambulatoire reste à intégrer.}
\end{frame}

\begin{frame}{Simulation dynamique : fermeture temporaire d'une salle}
  \small
  \centering
  \begin{tikzpicture}[x=1.0cm,y=1cm,>=Latex]
    \draw[very thick,draw=McmInk] (0,0) -- (10.3,0);
    \draw[line width=4pt,draw=McmAmber] (2.3,0) -- (5.5,0);
    \foreach \x/\lab in {0/08h,2.3/10h,5.5/12h,10.3/17h} {
      \draw[draw=McmInk] (\x,-0.12) -- (\x,0.12);
      \node[below=5pt,font=\scriptsize] at (\x,0) {\lab};
    }
    \node[above=8pt,text=McmAmber,font=\small\bfseries] at (3.9,0)
      {Salle 1 : aucune nouvelle entrée};
  \end{tikzpicture}
  \vspace{0.45cm}
  \begin{columns}[T,onlytextwidth]
    \begin{column}{0.49\textwidth}
      \begin{block}{Plan figé}
        \small Affectations et ordre initiaux conservés.
        Les cas suivants attendent que la salle soit disponible.
      \end{block}
    \end{column}
    \begin{column}{0.49\textwidth}
      \begin{exampleblock}{Plan réactif}
        \small Aux changements 10h/12h, le coordinateur soumet un
        nouveau planning pour les \textbf{cas en attente}.
      \end{exampleblock}
    \end{column}
  \end{columns}
  \vspace{0.10cm}
  \centering\small Intervention déjà commencée : elle se termine normalement.
  Proposition invalide : ancien plan conservé, garde d'exécution active.
  \source{Urgences, annulations et indisponibilités de lits ne sont pas encore des aléas exécutés dans ce simulateur.}
\end{frame}

\begin{frame}{Simulations : du planning à l'exécution}
  \footnotesize
\textbf{Une question commune : utiliser le bloc sans saturer les lits.}
\par\vspace{1.5mm}
\centering
\renewcommand{\arraystretch}{1.18}
\setlength{\tabcolsep}{5pt}
\begin{tabular}{@{}lp{3.15cm}p{6.25cm}@{}}
\toprule
\textbf{Évaluation} & \textbf{Population} & \textbf{Protocole}\\\midrule
Recherche statique & 300 patients / 80 vacations & 20 graines ; même budget de 3 s par méthode\\
Échelle historique & 582 patients / 280 vacations & 40 jours ; 4 s par méthode\\
Exécution simulée & 279 interventions / 22 jours & 2 salles, 08h--17h ; 42 ou 10 lits synthétiques\\
Incertitude LOS & 420 patients / 30 jours & 42 lits ; 10 graines ; 200 scénarios par planning\\\bottomrule
\end{tabular}
\vspace{0.05cm}
\begin{columns}[T,onlytextwidth]
\begin{column}{.49\textwidth}
\begin{block}{Évaluation statique}
Comparer la qualité des plannings à ressources et budget identiques.
\end{block}
\end{column}
\begin{column}{.49\textwidth}
\begin{exampleblock}{Exécution simulée}
Observer les interventions terminées, les dépassements et l'occupation des lits.
\end{exampleblock}
\end{column}
\end{columns}
\end{frame}

\begin{frame}{Scénarios : prévoir, puis confronter au réel}
  \footnotesize
\begin{block}{Du scénario à l'évaluation}
\centering
\textbf{Durées disponibles} (salle et séjour estimés)
$\longrightarrow$
\textbf{Planning} (bloc et lits)
$\longrightarrow$
\textbf{Évaluation} (exécution ou scénarios LOS)
\end{block}
\vspace{0.08cm}
\begin{columns}[T,onlytextwidth]
\begin{column}{.48\textwidth}
\begin{block}{Exécution historique : simulation SMA}
\textbf{Mêmes arrivées et durées observées.}\\
Médianes, prédictions ML ou oracle.\\
Salle 1 fermée de \textbf{10h à 12h} ; plan figé ou replanification des cas en attente.\\
Les actes commencés se terminent normalement.
\end{block}
\end{column}
\begin{column}{.48\textwidth}
\begin{block}{Séjours incertains : Monte Carlo LOS}
Planifier avec la prédiction ponctuelle ou la \textbf{borne haute conformelle}.\\
Planning figé ; \textbf{200 scénarios}, tirages communs.\\
Perturbation uniforme $\pm1{,}02$ jour, arrondie ; plancher d'un jour.\\
Un séjour prolongé mobilise davantage de lits.
\end{block}
\end{column}
\end{columns}
\end{frame}